\documentclass[11pt]{book}
\usepackage{amsmath,amssymb,amsthm}
\usepackage[paperwidth=170mm,paperheight=240mm,margin=18mm]{geometry}
\newtheorem{theorem}{Theorem}[section]
\newtheorem{lemma}{Lemma}[section]
\newtheorem{definition}{Definition}[section]
\begin{document}

\section{The Probabilistic Method}

\begin{definition}
A graph \(G=(V,E)\) is \emph{sparse} if \(|E| \le 3|V|-6\).
\end{definition}

\begin{theorem}[Tur\'an]
Every graph on \(n\) vertices with more than \(\frac{n^2}{4}\) edges contains a triangle.
\end{theorem}

\begin{proof}
Let \(G\) be triangle-free with \(n\) vertices and maximum number of edges.
For every edge \(uv\), the neighborhoods satisfy \(N(u)\cap N(v)=\varnothing\),
hence \(d(u)+d(v)\le n\). Summing over all edges gives
\[
\sum_{uv\in E}\bigl(d(u)+d(v)\bigr)=\sum_{v\in V} d(v)^2 \le n|E|.
\]
By Cauchy--Schwarz, \(\sum_v d(v)^2 \ge (2|E|)^2/n\), so \((2|E|)^2 \le n^2|E|\),
which yields \(|E|\le n^2/4\). \qed
\end{proof}

\begin{lemma}
If a graph has \(n\) vertices and \(e\) edges, then its average degree is \(2e/n\).
\end{lemma}

The celebrated Erd\H{o}s--R\'enyi bound shows that the random graph
\(G(n,p)\) with \(p=\frac{\log n}{n}\) is almost surely connected.

\end{document}
